\section{Gaussian squares: two scalar mechanisms}\label{sec:gaussian}


Now take $W=G^2$, where $G$ is standard Gaussian, and put
$X_{\mathrm{sq}}=X_{G^2}$.  Here $m_2=3$.  In particular, the product and
jump descriptions above give the same closed subspace with the same norm.
We first establish the Gaussian difference calculation that both estimates
will use, while keeping their subsequent probability and convexity steps
separate.

For fixed real coefficients $a_1,\ldots,a_k$ and a fixed $b\in\R$, put
\[
 H=b+\sum_i a_i(G_i^2-1),\qquad
 s=\Big(\sum_i a_i^2\Big)^{1/2},
\]
where the $G_i$ are independent standard Gaussians.  Let $H'$ be an
independent copy with the same coefficients and shift.  The Gaussian polarization identity used in decoupling
\cite[proof of Lemma~3.2]{LevinaVershynin2012} is especially simple here.
The orthogonal
change of variables
$A_i=(G_i+G_i')/\sqrt2$, $C_i=(G_i-G_i')/\sqrt2$ shows that
\[
 H-H'\ \text{has law}\ 2\sum_i a_iA_iC_i,
\]
where all $A_i,C_i$ are independent standard Gaussians.  Conditionally on
the $A_i$, this is Gaussian with variance $4V$, where
\[
 V=\sum_i a_i^2A_i^2,\qquad
 \E V=s^2,\qquad
 \E V^2=s^4+2\sum_i a_i^4\le3s^4.
\]
This is the elementary second-moment step usually called the
Paley--Zygmund inequality \cite{PaleyZygmund1932Series,PaleyZygmund1932Note}.
We include the argument: for $s>0$, the inequality
$\E[V\mathbf1_{\{V\ge s^2/2\}}]\ge s^2/2$ and Cauchy--Schwarz yield
$\PP(V\ge s^2/2)\ge1/12$.  On this event the conditional standard
deviation of $H-H'$ is at least $\sqrt2s$.  Consequently, for every $K>0$,
\begin{equation}\label{prob:gaussian-difference}
 \PP(|H-H'|\ge2Ks)
 \ge q_K,\qquad q_K=\frac{\PP(|G|\ge\sqrt2K)}{12}>0.
\end{equation}
The event in~\eqref{prob:gaussian-difference} has probability one when
$s=0$, so the same statement remains valid in that case.

\begin{proposition}[Shifted Gaussian tails]\label{prob:gaussian-24}
With $C=2+\sqrt2$, for every $t,\lambda>0$,
\begin{equation}\label{eq:gaussian-parameter-24}
 \whdelta_{X_{\mathrm{sq}}}(t)
 \ge\frac{\PP(|G|\ge\sqrt2\lambda)}{24}
       \left(\frac{t\lambda}C-1\right)_+.
\end{equation}
\end{proposition}
\begin{proof}
For the preceding fixed-coefficient variables, $|H-H'|\ge2\lambda s$
requires $|H|\ge\lambda s$ or $|H'|\ge\lambda s$.  Their equal laws and
\eqref{prob:gaussian-difference} therefore imply the shift-uniform estimate
\begin{equation}\label{prob:gaussian-shifted}
 \PP(|H|\ge\lambda s)\ge p_\lambda,
 \qquad p_\lambda=\frac{\PP(|G|\ge\sqrt2\lambda)}{24}.
\end{equation}
Again the zero-coefficient case is immediate.

For a finite-product $y\in F_T$, condition on the old coordinates and all
$B_i$ in~\eqref{prob:first-exit}.  It has the required form with
$a_i=B_i$ and shift $b=d+\sum_iB_i$.  Thus, for $x_0\in E_T$,
\begin{align}
 M_t(x_0,y)
 &\ge p_\lambda\E(t\lambda\sigma-|x_0|)_+\notag\\
 &\ge p_\lambda\bigl(t\lambda\E\sigma-\norm{x_0}\bigr)
 \ge p_\lambda\left(\frac{t\lambda}{C}\norm{y}-\norm{x_0}\right).
 \label{prob:gaussian-24-gain}
\end{align}
The last step uses~\eqref{prob:first-moment-centered} with $m_2=3$.
Lemma~\ref{lem:affine} now proves~\eqref{eq:gaussian-parameter-24}.
Taking $\lambda=4C/t$ gives the first Gaussian bound in
Theorem~\ref{thm:main} for every $t>0$.
\end{proof}

\begin{proposition}[Convex comparison with a Gaussian difference]\label{prob:gaussian-12}
For every $t,K>0$, the same norm satisfies
\begin{equation}\label{eq:gaussian-parameter-12}
 \whdelta_{X_{\mathrm{sq}}}(t)
 \ge\frac{\PP(|G|\ge\sqrt2K)}{12}
       \left(\frac{tK}{2\sqrt3}-1\right)_+.
\end{equation}
\end{proposition}
\begin{proof}
For a finite-product $y\in F_T$, fix the old coordinates and the $B_i$,
and resample only the first-exit Gaussian squares to obtain $\widetilde y$.
These variables have the same conditional law.  Applying
\eqref{prob:gaussian-difference} to their difference gives
\begin{equation}\label{prob:conditional-gaussian-difference}
 \PP\bigl(|y-\widetilde y|\ge2K\sigma
       \mid\mathcal G_T,(B_i)_i\bigr)\ge q_K.
\end{equation}
For fixed old coordinates and $x_0\in E_T$, the function
$\phi(v)=(t|v|-|x_0|)_+$ is even and convex.  Hence
\[
 \phi\bigl((y-\widetilde y)/2\bigr)
 \le\tfrac12\bigl(\phi(y)+\phi(-\widetilde y)\bigr).
\]
Taking conditional expectations given $\mathcal H_T$ and then integrating gives
$\E\phi(y)\ge\E\phi((y-\widetilde y)/2)$.
Combining this inequality with~\eqref{prob:conditional-gaussian-difference}
and~\eqref{prob:first-moment-signed}, now with $m_2=3$, yields
\begin{align}
 M_t(x_0,y)
 &\ge\E(t|y-\widetilde y|/2-|x_0|)_+\notag\\
 &\ge q_K\E(tK\sigma-|x_0|)_+\notag\\
 &\ge q_K\left(\frac{tK}{2\sqrt3}\norm{y}-\norm{x_0}\right).
 \label{prob:gaussian-12-gain}
\end{align}
This argument uses the probability of a difference directly; the convexity
step is what preserves the denominator $12$ in $q_K$.
Apply Lemma~\ref{lem:affine} to this affine bound. This proves
\eqref{eq:gaussian-parameter-12}. Taking $K=8\sqrt3/t$ yields the second
Gaussian bound in Theorem~\ref{thm:main} for every $t>0$.
\end{proof}

The two proofs control different scalar quantities. Proposition~\ref{prob:gaussian-24}
first estimates a shifted original variable and pays a factor of two in
the union bound. Proposition~\ref{prob:gaussian-12} estimates a difference
and uses even convexity directly in the midpoint increment, preserving
the denominator $12$. Their first-moment constants are respectively
$2+\sqrt2$ and $2\sqrt3$. Neither proof requires the original shifted
variable to be symmetric.
